The verification component evaluates execution traces against safety rules expressed in Linear Temporal Logic over finite traces (LTLf). Since task executions terminate after a finite number of steps, LTLf provides a more appropriate semantics than classical LTL, whose formulas are interpreted over infinite executions. The framework uses the \texttt{flloat} library to parse LTLf formulas and evaluate their satisfaction over the collected traces.

Before verification, the translated execution trace is converted into the Boolean interpretation required by \texttt{flloat}. Each execution step is represented as a dictionary mapping grounded predicates to Boolean values, while the corresponding execution predicate (e.g., \texttt{exec\_\_PickupObject\_\_knife\_1}) is merged into the same representation. Parameterized LTL templates generated by the rule generation module are then instantiated with the concrete object identifiers appearing in the current task, producing fully grounded formulas that can be evaluated directly over the trace.

A safety violation is defined as the first execution prefix for which an LTLf formula evaluates to \texttt{False}. After parsing each formula, the checker evaluates its truth over the complete trace. If a formula is violated, the earliest violating step is identified by evaluating progressively longer prefixes of the trace until the formula becomes unsatisfied. The checker reports the violated rule, the grounded formula, the action executed at the violating step, and the corresponding step index.

Unlike traditional offline verification approaches, the proposed framework also supports runtime monitoring. Rather than waiting for the task to terminate, every newly generated execution step is immediately translated and appended to an accumulated trace. After each action, the monitor invokes the LTL checker on the current execution prefix. If any safety rule is violated, a \texttt{SafetyViolation} exception is raised, the violation is recorded in a log file together with the violated rules and timestamp, and the agent is immediately interrupted. Consequently, unsafe behaviours are prevented from continuing beyond the first detected violation, enabling online enforcement of temporal safety specifications rather than merely post-hoc analysis.